\section{FORMAL RESTATEMENT}
\textbf{Erd\H{o}s \#511; reconstruction of the negative answer, 2026-09-24.}
For a monic polynomial $f$, let $E(f)=\{z\in\mathbb C:|f(z)|<1\}$. Is the number of components of $E(f)$ of diameter greater than a fixed $c>1$ bounded independently of the degree?

\section{QUICK LITERATURE AND CONTEXT CHECK}
Huang's \emph{Many Lemniscates with large diameter}, version 2, gives counterexamples for every $1<c<4$ and attributes the earlier result to Pommerenke. We reconstruct the capacity argument, taking care that the question uses an open lemniscate.

\section{OBJECTIVE AND ATTACK PLAN}
Place many disjoint long rectangles in a thin ellipse of capacity one, approximate their union by a polynomial lemniscate, and make the polynomial monic by a dilation that cannot shrink distances.

\section{WORK}
We use three classical potential-theoretic inputs explicitly. The capacity of a full compact set with an exterior conformal map $az+O(1)$ is $|a|$. Capacity is monotone under inclusion. For a polynomial $q$ of degree $d$ and leading coefficient $a$, the compact lemniscate $\{|q|\leq1\}$ has capacity $|a|^{-1/d}$. We also use Hilbert's lemniscate approximation theorem: for a compact set $K$ with connected complement and every neighborhood $U$ of $K$, there is a polynomial $q$ with
\[
 K\subseteq\{|q|\leq1\}\subseteq U.
\]
These are the external inputs; the construction and all normalizations follow below.

Fix $1<c<4$ and a desired number $N$. Choose $0<s<1$ and $d_0$ with
$c<d_0<2(1+s)$. The map $z\mapsto z+s/z$ conformally maps the exterior unit disk onto the exterior of the ellipse with semiaxes $1+s$ and $1-s$. Its leading coefficient is one, so the filled ellipse $Q$ has capacity one. Because its horizontal cross-sections near the center have length greater than $d_0$, its interior contains $N$ pairwise disjoint thin closed horizontal rectangles, each of length $d_0$. Let their union be $K$. The complement of this finite union of disjoint rectangles is connected.
Choose a neighborhood $U$ consisting of disjoint neighborhoods of the rectangles, with closure in the interior of $Q$. Apply the approximation theorem with a still smaller neighborhood whose closure is in $U$, obtaining $q$ as above. There exists $\eta>0$ such that
\[
 L:=\{|q|\leq1+\eta\}\subset U.
\]
Indeed, on the complement of $U$ in a sufficiently large closed disk, $|q|$ has minimum strictly greater than one, and outside that disk it is greater than two. Choose $\eta$ smaller than these margins. Then $q_0=q/(1+\eta)$ satisfies
\[
 K\subset\{|q_0|<1\},\qquad \{|q_0|\leq1\}\subset Q.
\]
Every rectangle, being connected, lies in one component of the open lemniscate. Different rectangles lie in different components since these components stay in the disjoint pieces of $U$. Each of the selected $N$ components has diameter at least $d_0$.

If $a$ is the leading coefficient of $q_0$ and $d=\deg q_0$, monotonicity and the capacity formula give $|a|^{-1/d}\leq1$, hence $|a|\geq1$. Choose $w$ with $w^d=a$ and set $f(z)=q_0(z/w)$. This polynomial is monic, and $E(f)=wE(q_0)$. Thus the dilation multiplies diameters by $|w|\geq1$. It preserves the number of components and leaves at least $N$ components with diameter greater than $c$.

\section{RESULT AND LIMITATIONS}
\textbf{Attempt status: solved.} For each $1<c<4$ there is no degree-independent bound. This suffices to disprove the proposed assertion. The reconstruction depends on the explicitly stated classical capacity and lemniscate approximation theorems; their proofs and a classification at the endpoint $c=4$ are outside this attempt.

\section{SOURCES}
\url{https://www.erdosproblems.com/511}; W. Huang, \emph{Many Lemniscates with large diameter}, version 2, \url{https://arxiv.org/html/2509.11597v2}.

\textbf{Completion rate estimate: 100\%.}
